\chapter{Addressing Complex Engineering Problems and Activities}

\section{Introduction}
The International Engineering Alliance's graduate-attribute
framework~\cite{iea2021gapc} defines the
attributes of a Complex Engineering Problem (WP1 to WP7) and of Complex
Engineering Activities (EA1 to EA5). A problem qualifies as complex when it
demands in-depth knowledge (WP1) together with some or all of the remaining
attributes; an activity qualifies when its execution shows the corresponding
range of resources, interactions, and innovation. This chapter maps the
thesis onto those attributes, and each claim points at the section, table,
or figure of this report where the evidence sits.

\section{Complex Engineering Problems Associated with the Work}

Seven attributes mark a problem as one of complex engineering. Each is taken in
turn below and matched against what this thesis actually required, from the depth
of background knowledge needed to the conflicting requirements that had to be
traded off. Every claim points at the section, table, or figure that carries its
evidence, so that the mapping can be checked rather than taken on trust.

\begin{itemize}
  \item \textbf{WP1: Depth of knowledge required.} The default detector was
        inadequate for the sub-ten-pixel regime measured in this dataset. The
        investigation required working knowledge of one-stage detection
        architectures at the layer level (Table~\ref{tab:layers}), of
        channel and spatial attention (Section~\ref{sec:cbam3}), of selective
        state-space models and their scan mechanics (Section~\ref{sec:mamba}), and of
        the loss geometry that makes few-pixel boxes unstable under strict
        IoU (Sections~\ref{sec:loss3} and~\ref{sec:metrics}), a fundamentals-first analysis rather
        than tool use.
  \item \textbf{WP2: Range of conflicting requirements.} The core tension is
        technical and non-technical at once: recall on the smallest targets
        against a parameter and latency budget for airborne hardware, and
        the humanitarian value of the system against its surveillance
        potential (Chapter~V). The resolution is measured, not asserted: the
        P2 scale buys very-tiny recall at $+0.5$~M parameters and $+1$~ms
        (Table~\ref{tab:main}), and the recall-weighted $F_2$ metric encodes
        the ethical asymmetry directly into the evaluation.
  \item \textbf{WP3: Depth of analysis, no obvious solution.} Which of
        attention, resolution, or sequence modelling helps tiny-target
        detection had no a-priori answer; each had a plausible mechanism. The
        additive ablation (Section~\ref{sec:protocol}) provides the comparison,
        and its result distinguishes one useful scale change, one efficient
        substitution, and one costly null result.
  \item \textbf{WP4: Familiarity of issues.} The dominant issues are
        infrequently encountered in standard detection practice: targets
        smaller than one detection cell (Figure~\ref{fig:strideprob}), a
        dataset whose instances are $99.6\%$ below 32~px
        (Figure~\ref{fig:sizedist}), and a training collapse specific to
        the state-space blocks at epoch 28 (Figure~\ref{fig:training}).
  \item \textbf{WP5: Beyond standards and codes.} No engineering standard
        prescribes how to evaluate or deploy a tiny-human detector for
        search and rescue. The evaluation protocol, the frozen-split and
        provenance rules, and the operational threshold policy
        (Sections~\ref{sec:metrics} and~\ref{sec:protocol}) had to be designed for the problem.
  \item \textbf{WP6: Stakeholder scope.} The stakeholders of a rescue
        detector, operators who act on its boxes, agencies bound by aviation
        and data-protection rules, and the people being imaged, have needs
        that pull in different directions; Chapter~V works through those
        obligations and the design choices that answer them.
  \item \textbf{WP7: Interdependence.} The system is a chain of coupled
        sub-problems: dataset properties drive the architecture
        (Section~\ref{sec:p2novel}), the architecture drives memory behaviour and hence
        the training configuration (Section~\ref{sec:setup}), and the evaluation design
        determines whether any of it can be attributed. The single-variable
        ablation exists precisely to cut through that interdependence.
\end{itemize}

\section{Complex Engineering Activities Associated with the Work}

Where the previous section described the nature of the problem, this one
describes what carrying out the work involved. Five attributes define a complex
engineering activity, and each is set against the resources the project drew on,
the way the experiments were organised, and the judgements that had to be made
with incomplete information. The same referencing convention applies.

\begin{itemize}
  \item \textbf{EA1: Range of resources.} The work managed diverse
        resources: a 10{,}215-image benchmark with frozen splits, pretrained
        weights, a single 16~GB consumer GPU shared across 54 hours of
        training (Table~\ref{tab:carbon}), open-source frameworks, and a
        metric pipeline producing per-run manifests, CSV data, and plots so
        that every reported number is regenerable.
  \item \textbf{EA2: Level of interactions.} Theoretical design met hardware
        limits directly. The stride-4 scale quadruples head activation
        memory, which forced the physical batch to 8 with gradient
        accumulation holding the effective batch at 16 (Section~\ref{sec:setup}), and
        the automatic optimiser choice diverged on that same architecture,
        which forced the pinned AdamW configuration; resolving these
        interactions without breaking comparability was the day-to-day
        engineering of the thesis.
  \item \textbf{EA3: Innovation.} The individual modules are established; the
        contribution lies in how they are tested and selected. Attention is
        substituted rather than appended, the state-space blocks are isolated
        in the neck of an otherwise unchanged detector, and tiny-object recall
        and $F_2$ are treated as decision metrics. The later assignment blend
        then modifies positive-sample selection without adding inference cost
        (Section~\ref{sec:assignmethod}).
  \item \textbf{EA4: Consequences for society and the environment.} The
        activity was run with its externalities measured: a carbon tracker
        and power sampler on every run (9.0~kg CO$_2$ total,
        Table~\ref{tab:carbon}), an ethical position on dual use and
        training-data dignity (Chapter~V), and a deliverable whose small
        size lowers the energy cost of every future deployment.
  \item \textbf{EA5: Familiarity, extension beyond prior experience.} The
        project extended well past coursework: reading and adapting current
        detection, attention, and state-space literature; diagnosing a
        silent model-rebuild fault that had invalidated an earlier line of
        results; and building failure-tolerant training infrastructure
        (smoke tests, resume, out-of-memory ladders) for a load-shedding
        environment, none of which is met in standard undergraduate
        laboratory work. The extended phase (Section~\ref{sec:partb}) added
        two more activities of this kind: checkpoint-resumable training
        chains run unattended on shared laboratory machines, and a
        data-engineering pipeline for the real-imagery sets, from frame
        extraction and perceptual-hash de-duplication through assisted
        annotation and automated quality audits
        (Section~\ref{sec:adaptmethod}).
\end{itemize}

\section{Conclusion}
Measured against the Washington Accord attributes~\cite{iea2021gapc}, the
thesis presents a complex engineering problem, in-depth knowledge applied to
an unfamiliar, standard-less task with conflicting requirements and coupled
sub-problems, and its execution shows the corresponding complex activities,
from resource and interaction management to innovation and measured
consequences. Each attribute above is tied to a numbered table, figure, or
section so that the mapping can be checked against the work reported here.
